% =====================================================================
%  PRÉAMBULE — « ANGLAIS : LE GRAND LIVRE D'AUTO-FORMATION »
%  Méthode complète pour francophones — 12 mois, A1 -> B2/C1
%
%  > Moteur conseillé : pdfLaTeX (2 passes pour la table des matières)
%  > Fonctionne sur Overleaf (TeX Live complet) et en local
%    (TeX Live / MiKTeX). Aucun fichier image externe n'est requis.
%  > Si un paquet manque, MiKTeX le télécharge automatiquement ;
%    sur TeX Live :  sudo tlmgr install tcolorbox titlesec pifont
% =====================================================================

% ---------- Encodage et polices -------------------------------------
\usepackage[T1]{fontenc}
\usepackage[utf8]{inputenc}
\usepackage{lmodern}
\usepackage{microtype}
\usepackage{textcomp}
\usepackage{pifont}

% ---------- Langues : français (principal) + anglais ----------------
\usepackage[english,french]{babel}
% Typographie française correcte ; l'anglais est inséré via \en{} ou
% \selectlanguage{english} pour les longs passages.
\frenchbsetup{
  AutoSpacePunctuation = true,
  StandardItemizeEnv   = true,
  StandardEnumerateEnv = true,
  ThinSpaceInFrenchNumbers = true,
  FrenchFootnotes      = false,
  IndentFirst          = false
}

% ---------- Mise en page --------------------------------------------
\usepackage[a4paper,top=2.2cm,bottom=2.2cm,left=2.1cm,right=2.1cm,
            headheight=14pt,footskip=1.1cm]{geometry}
\usepackage{fancyhdr}
\usepackage{titlesec}
\usepackage{multicol}
\usepackage{longtable}
\usepackage{array}
\usepackage{tabularx}
\usepackage{colortbl}
\usepackage{ragged2e}
\usepackage[usenames,dvipsnames,table]{xcolor}

% ---------- Couleurs du livre ---------------------------------------
\definecolor{bleuNuit}{HTML}{0F3D6E}
\definecolor{bleuClair}{HTML}{E9F1FA}
\definecolor{bleuMoyen}{HTML}{3D7EBF}
\definecolor{vert}{HTML}{1B7A43}
\definecolor{vertClair}{HTML}{E8F6EE}
\definecolor{rouge}{HTML}{AE2B2B}
\definecolor{rougeClair}{HTML}{FCECEC}
\definecolor{orange}{HTML}{C4720A}
\definecolor{orangeClair}{HTML}{FFF3E1}
\definecolor{violet}{HTML}{5B3A8E}
\definecolor{violetClair}{HTML}{F2ECFA}
\definecolor{gris}{HTML}{5A5A5A}
\definecolor{grisClair}{HTML}{F3F3F3}
\definecolor{grisTresClair}{HTML}{FAFAFA}
\definecolor{tabHead}{HTML}{0F3D6E}
\definecolor{tabAlt}{HTML}{F4F7FB}

% ---------- Symboles utilitaires ------------------------------------
\newcommand{\ok}{\ding{51}\,}
\newcommand{\ko}{\ding{55}\,}
\newcommand{\etoile}{\ding{72}\,}
\newcommand{\ra}{\textrightarrow\,}
\newcommand{\Ra}{\ensuremath{\Rightarrow}\,}
\newcommand{\cible}{\ding{119}\,}

% ---------- Phonétique : IPA (si tipa est installé) ------------------
\newif\ifavecTIPA
\IfFileExists{tipa.sty}{\avecTIPAtrue}{\avecTIPAfalse}
%  Voyelles :   \ipaee iː  \ipaiot ɪ  \ipaE ɛ/e  \ipae æ  \ipaa ɑː  \ipao ɔː
%               \ipaups ʊ  \ipaoo uː  \ipalv ʌ  \ipaer ɜː  \ipeschwa ə
%  Diphtongues: \ipaei eɪ  \ipaai aɪ  \ipaOi ɔɪ  \ipaau aʊ  \ipaou əʊ/oʊ
%  Consonnes : \ipatheta θ  \ipadh ð  \ipash ʃ  \ipazh ʒ  \ipang ŋ
%  Divers   :  \ipalong ː   \ipa{...} séquence libre (notation TIPA)
\ifavecTIPA
  \usepackage[safe]{tipa}
  \newcommand{\ipash}{\textipa{S}}        % ʃ
  \newcommand{\ipazh}{\textipa{Z}}        % ʒ
  \newcommand{\ipang}{\textipa{N}}        % ŋ
  \newcommand{\ipeschwa}{\textipa{@}}     % ə
  \newcommand{\ipatheta}{\textipa{T}}     % θ
  \newcommand{\ipadh}{\textipa{D}}        % ð
  \newcommand{\ipae}{\ae}                 % æ
  \newcommand{\ipaE}{\textipa{E}}         % ɛ
  \newcommand{\ipaa}{\textipa{A}}         % ɑ
  \newcommand{\ipao}{\textipa{O}}         % ɔ
  \newcommand{\ipaer}{\textipa{3}}        % ɜ
  \newcommand{\ipaups}{\textipa{U}}       % ʊ
  \newcommand{\ipaiot}{\textipa{I}}       % ɪ
  \newcommand{\ipalv}{\textipa{V}}        % ʌ
  \newcommand{\ipaoo}{\textipa{u}}        % u
  \newcommand{\ipaee}{\textipa{i}}        % i
  \newcommand{\ipalong}{\textipa{:}}      % ː
  \newcommand{\ipaei}{\textipa{eI}}       % eɪ
  \newcommand{\ipaai}{\textipa{aI}}       % aɪ
  \newcommand{\ipaOi}{\textipa{OI}}       % ɔɪ
  \newcommand{\ipaau}{\textipa{aU}}       % aʊ
  \newcommand{\ipaou}{\textipa{@U}}       % əʊ
  \newcommand{\ipaoou}{\textipa{oU}}      % oʊ (AmE)
  \newcommand{\ipa}[1]{\textipa{#1}}
\else
  % ---- Repli : transcription lisible si le paquet tipa est absent ----
  \newcommand{\ipash}{ch}
  \newcommand{\ipazh}{j}
  \newcommand{\ipang}{ng}
  \newcommand{\ipeschwa}{\textsf{@}}
  \newcommand{\ipatheta}{\ensuremath{\theta}}
  \newcommand{\ipadh}{dh}
  \newcommand{\ipae}{\ae}
  \newcommand{\ipaE}{\`e}
  \newcommand{\ipaa}{a:}
  \newcommand{\ipao}{o:}
  \newcommand{\ipaer}{eu:}
  \newcommand{\ipaups}{ou}
  \newcommand{\ipaiot}{i}
  \newcommand{\ipalv}{e}
  \newcommand{\ipaoo}{ou:}
  \newcommand{\ipaee}{i:}
  \newcommand{\ipalong}{:}
  \newcommand{\ipaei}{\'e-i}
  \newcommand{\ipaai}{a-i}
  \newcommand{\ipaOi}{o-i}
  \newcommand{\ipaau}{a-ou}
  \newcommand{\ipaou}{o-ou}
  \newcommand{\ipaoou}{o-ou}
  \newcommand{\ipa}[1]{#1}
\fi

% Prononciation « à la française » : toujours disponible, sans paquet
\newcommand{\pron}[1]{{\color{violet}\sffamily\small[#1]}}

% ---------- Anglais en ligne et blocs bilingues ---------------------
% \en{}  : passage en anglais (typographie anglaise, pas d'espace avant ! ? :)
% \enb{} : idem en gras   \eni{} : idem en italique
\newcommand{\en}[1]{\foreignlanguage{english}{#1}}
\newcommand{\enb}[1]{\foreignlanguage{english}{\textbf{#1}}}
\newcommand{\eni}[1]{\foreignlanguage{english}{\textit{#1}}}

% Exemple bilingue : anglais + traduction française en retrait
\newcommand{\ex}[2]{%
  \par\smallskip
  {\selectlanguage{english}#1\par}%
  \hspace*{1.2em}{\footnotesize\color{gris}\selectlanguage{french}\itshape #2}\par}

% Phrase utile (tableaux) : anglais | français
\newcommand{\phr}[2]{\en{#1} & \selectlanguage{french}#2 \tabularnewline}

% Mot de vocabulaire : mot | prononciation française | sens
\newcommand{\mot}[3]{\textbf{\en{#1}} & \pron{#2} & \selectlanguage{french}#3 \tabularnewline}

% Ligne de dialogue : personnage, réplique anglaise, traduction
\newcommand{\dial}[3]{%
  \par\smallskip
  {\bfseries\color{bleuNuit}\selectlanguage{french}#1}\enspace\textcolor{orange}{\textbullet}\enspace
  {\selectlanguage{english}#2\par}%
  \hspace*{1.6em}{\footnotesize\color{gris}\selectlanguage{french}\itshape #3}\par}

% Réponse courte / variante dans un dialogue
% Flèches d'intonation : sûres en mode texte comme en mode maths
\newcommand{\monte}{\ensuremath{\nearrow}}
\newcommand{\descend}{\ensuremath{\searrow}}

\newcommand{\variante}[1]{\par\hspace*{1.6em}{\footnotesize\selectlanguage{english}\textcolor{bleuMoyen}{#1}}\par}

% Trou à compléter dans les exercices
\newcommand{\trou}[1][2.2cm]{\rule{#1}{0.4pt}}
\newcommand{\troul}[1]{\underline{\hspace{#1}}}

% ---------- Boîtes colorées (tcolorbox si disponible) ----------------
\IfFileExists{tcolorbox.sty}{%
  \usepackage[skins,breakable]{tcolorbox}
  \tcbset{
    boite/.style={enhanced jigsaw, breakable, boxrule=0pt,
      leftrule=3pt, arc=0pt, outer arc=0pt, left=7pt, right=7pt,
      top=6pt, bottom=6pt, toptitle=4pt, bottomtitle=4pt,
      fonttitle=\bfseries\selectlanguage{french},
      coltitle=white, fontupper=\small, fontlower=\small},
    sansTitre/.style={enhanced jigsaw, breakable, boxrule=0.5pt, arc=2pt,
      left=6pt, right=6pt, top=5pt, bottom=5pt, fontupper=\small}
  }
  \newtcolorbox{regle}[1][]{boite, colback=bleuClair, colframe=bleuClair,
    borderline west={3pt}{0pt}{bleuNuit}, coltitle=white,
    title={\selectlanguage{french}\ok RÈGLE #1}, colbacktitle=bleuNuit}
  \newtcolorbox{piege}[1][]{boite, colback=rougeClair, colframe=rougeClair,
    borderline west={3pt}{0pt}{rouge},
    title={\selectlanguage{french}\ko PIÈGE FRANCOPHONE #1}, colbacktitle=rouge}
  \newtcolorbox{astuce}[1][]{boite, colback=vertClair, colframe=vertClair,
    borderline west={3pt}{0pt}{vert},
    title={\selectlanguage{french}\etoile ASTUCE #1}, colbacktitle=vert}
  \newtcolorbox{mission}[1][]{boite, colback=orangeClair, colframe=orangeClair,
    borderline west={3pt}{0pt}{orange},
    title={\selectlanguage{french}\cible MISSION DU JOUR #1}, colbacktitle=orange}
  \newtcolorbox{memo}[1][]{boite, colback=violetClair, colframe=violetClair,
    borderline west={3pt}{0pt}{violet},
    title={\selectlanguage{french}\ding{46}\ À RETENIR #1}, colbacktitle=violet}
  \newtcolorbox{dialoguebox}[1][]{sansTitre, colback=grisTresClair, colframe=bleuMoyen,
    title={\selectlanguage{french}\bfseries DIALOGUE #1}, coltitle=white,
    colbacktitle=bleuMoyen}
  \newtcolorbox{exo}[1][]{sansTitre, colback=white, colframe=bleuNuit,
    title={\selectlanguage{french}\bfseries EXERCICE #1}, coltitle=white,
    colbacktitle=bleuNuit}
  \newtcolorbox{corrige}[1][]{sansTitre, colback=vertClair, colframe=vert,
    title={\selectlanguage{french}\bfseries CORRIGÉ #1}, coltitle=white,
    colbacktitle=vert}
  \newtcolorbox{tableau}[1][]{sansTitre, colback=white, colframe=gris,
    title={\selectlanguage{french}\bfseries #1}, coltitle=white,
    colbacktitle=gris}
}{%
  % -------- Repli sans tcolorbox : contenu toujours lisible ----------
  \newenvironment{regle}[1][]{\par\medskip\noindent\begin{minipage}{\linewidth}
      \color{bleuNuit}\bfseries\selectlanguage{french}\ok RÈGLE #1\par\normalcolor\small\selectlanguage{french}}
      {\end{minipage}\par\medskip}
  \newenvironment{piege}[1][]{\par\medskip\noindent\begin{minipage}{\linewidth}
      \color{rouge}\bfseries\selectlanguage{french}\ko PIÈGE FRANCOPHONE #1\par\normalcolor\small\selectlanguage{french}}
      {\end{minipage}\par\medskip}
  \newenvironment{astuce}[1][]{\par\medskip\noindent\begin{minipage}{\linewidth}
      \color{vert}\bfseries\selectlanguage{french}\etoile ASTUCE #1\par\normalcolor\small\selectlanguage{french}}
      {\end{minipage}\par\medskip}
  \newenvironment{mission}[1][]{\par\medskip\noindent\begin{minipage}{\linewidth}
      \color{orange}\bfseries\selectlanguage{french}\cible MISSION DU JOUR #1\par\normalcolor\small\selectlanguage{french}}
      {\end{minipage}\par\medskip}
  \newenvironment{memo}[1][]{\par\medskip\noindent\begin{minipage}{\linewidth}
      \color{violet}\bfseries\selectlanguage{french}\ding{46}\ À RETENIR #1\par\normalcolor\small\selectlanguage{french}}
      {\end{minipage}\par\medskip}
  \newenvironment{dialoguebox}[1][]{\par\medskip\noindent\begin{minipage}{\linewidth}
      \bfseries\selectlanguage{french}DIALOGUE #1\par\mdseries\small}
      {\end{minipage}\par\medskip}
  \newenvironment{exo}[1][]{\par\medskip\noindent\begin{minipage}{\linewidth}
      \bfseries\selectlanguage{french}EXERCICE #1\par\mdseries\small}
      {\end{minipage}\par\medskip}
  \newenvironment{corrige}[1][]{\par\medskip\noindent\begin{minipage}{\linewidth}
      \bfseries\selectlanguage{french}CORRIGÉ #1\par\mdseries\small}
      {\end{minipage}\par\medskip}
  \newenvironment{tableau}[1][]{\par\medskip\noindent\begin{minipage}{\linewidth}
      \bfseries\selectlanguage{french}#1\par\mdseries\small}
      {\end{minipage}\par\medskip}
}

% ---------- Environnements simples (sans paquet) ---------------------
% Bloc 100 % anglais (dialogues, textes, listes en anglais)
\newenvironment{anglais}{\selectlanguage{english}\par\smallskip}{\par\smallskip\selectlanguage{french}}
% Bloc de vocabulaire en trois colonnes
\newenvironment{tabvocab}
  {\par\smallskip\selectlanguage{french}\noindent
   \begin{longtable}{@{}p{4.1cm}p{3.1cm}p{8.4cm}@{}}
   \rowcolor{tabHead}\color{white}\bfseries Mot & \bfseries Prononciation & \bfseries Sens, exemple \\[2pt]
   \endfirsthead
   \rowcolor{tabHead}\color{white}\bfseries Mot & \bfseries Prononciation & \bfseries Sens, exemple \\[2pt]
   \endhead}
  {\end{longtable}\par\smallskip}

% ---------- Titres (parties, chapitres, sections) --------------------
\titleformat{\part}[display]
  {\normalfont\sffamily\color{bleuNuit}}
  {\filleft\huge\bfseries\partname\ \thepart}{12pt}
  {\filcenter\Huge\bfseries}
  [\vspace{6pt}{\color{orange}\titlerule[2pt]}]
\titleformat{\chapter}[hang]
  {\normalfont\sffamily\Huge\bfseries\color{bleuNuit}}
  {\thechapter\enspace\textcolor{orange}{\vrule width 4pt height 24pt}\enspace}{0pt}{\Huge}
\titlespacing*{\chapter}{0pt}{-18pt}{22pt}
\titleformat{\section}
  {\normalfont\sffamily\Large\bfseries\color{bleuNuit}}{\thesection}{0.8em}{}
\titlespacing*{\section}{0pt}{16pt}{8pt}
\titleformat{\subsection}
  {\normalfont\sffamily\large\bfseries\color{orange}}{\thesubsection}{0.8em}{}
\titlespacing*{\subsection}{0pt}{12pt}{6pt}
\titleformat{\subsubsection}
  {\normalfont\sffamily\normalsize\bfseries\color{gris}}{\thesubsubsection}{0.8em}{}
\titlespacing*{\subsubsection}{0pt}{10pt}{4pt}

% ---------- En-têtes et pieds de page --------------------------------
\pagestyle{fancy}
\fancyhf{}
\fancyhead[L]{\footnotesize\sffamily\color{gris}\nouppercase{\rightmark}}
\fancyhead[R]{\footnotesize\sffamily\color{bleuNuit}\thepage}
\fancyfoot[L]{\footnotesize\sffamily\color{gris}Anglais --- le grand livre d'auto-formation}
\renewcommand{\headrulewidth}{0.4pt}
\renewcommand{\headrule}{\hbox to\headwidth{\color{bleuMoyen}\leaders\hrule height \headrulewidth\hfill}}
\renewcommand{\footrulewidth}{0pt}
\fancypagestyle{plain}{%
  \fancyhf{}\fancyfoot[C]{\footnotesize\sffamily\thepage}%
  \renewcommand{\headrulewidth}{0pt}}

% ---------- Listes plus aérées ---------------------------------------
\usepackage{enumitem}
\setlist[itemize]{leftmargin=1.4em,itemsep=2pt,topsep=4pt,parsep=0pt}
\setlist[enumerate]{leftmargin=1.8em,itemsep=2pt,topsep=4pt,parsep=0pt}

% ---------- Tableaux : style par défaut ------------------------------
\setlength{\tabcolsep}{5pt}
\renewcommand{\arraystretch}{1.25}
\newcolumntype{L}[1]{>{\RaggedRight\arraybackslash}p{#1}}
\newcolumntype{C}[1]{>{\Centering\arraybackslash}p{#1}}
\newcommand{\entete}[1]{\rowcolor{tabHead}\color{white}\bfseries #1}

% ---------- Liens internes -------------------------------------------
\usepackage[colorlinks=true,linkcolor=bleuNuit,urlcolor=bleuMoyen,
            citecolor=vert,pdftitle={Anglais - le grand livre d'auto-formation},
            pdfauthor={},pdfsubject={Methode complete d'anglais pour francophones},
            bookmarksopen=true,bookmarksnumbered=true]{hyperref}

% ---------- Divers ---------------------------------------------------
\setlength{\parskip}{4pt}
\setlength{\parindent}{0pt}
\raggedbottom
\sloppy
\newcommand{\pagevierge}{\clearpage}
